\chapter{Results So Far}
\label{ch:results}

This chapter reports every experiment run up to 9 October 2026, roughly in the order they were made. All of them are on
the development split. Most are a single run on a handful of questions, made while the harness was being built, and
several were made to decide a design question rather than to measure the final system. They show where the harness
stands and what each component did when it was introduced; they do not yet answer the research question. Unless
stated otherwise, tables are produced by \code{eval/analysis/report.py} from the folder named in the caption, and the
model is \code{qwen3:8b}.

\section{First experiment: node degree on small graphs}

The first experiment asked 100 node-degree questions on the small development graphs (two per graph), with only two
tools available, \code{get_neighbors} and an edge counter, and no edge list in the prompt. It was planned as three
runs; the third was stopped after 12 questions. Over the 212 answers, accuracy was 99.5\% (95\% CI 98.6--100\%), at
about 1,356 tokens and 4.3 seconds per answer. The only failure was not a reasoning failure: the model called
\code{get_neighbors}, received the six neighbors, and then wrote its (correct) answer of 6 as a text tool call three
times, until loop detection stopped the run. This was the first sign that the final answer's format, not the
reasoning, would be the cheap model's weak point.
% source: results/harness_runs/degree_small_dev_v1, docs/learning_log.md (2026-10-04)

\section{Well-fitting tools}

Three times during the project a tool was added that answers a task directly, and each time the effect was the same.

\subsection{A direct edge test}

Table~\ref{tab:has-edge} compares five edge-existence questions on small graphs with and without \code{has_edge}.
Without it, the model spent 1,500--2,700 thinking tokens working out that \code{get_neighbors} answers the question
before making its first call; with it, 275--287 tokens every time.

\begin{table}[ht]
  \centering
  \small
  \caption[Ablation: the has\_edge tool]{Ablation of \code{has_edge}: 5 edge-existence questions, small development
  graphs, 1 run. Prompt tokens in these runs are inflated by the thinking-in-history bug (fixed later); accuracy and
  output tokens are unaffected.}
  \label{tab:has-edge}
  % source: results/harness_runs/has_edge_ablation_without, results/harness_runs/has_edge_ablation_with, docs/learning_log.md (2026-10-05)
  \begin{tabular}{lrr}
    \toprule
    & without \code{has_edge} & with \code{has_edge} \\
    \midrule
    correct & 5/5 & 5/5 \\
    output tokens (all 5 questions, mostly thinking) & 2,256 & 282 \\
    prompt tokens (all 5 questions) & 4,200 & 1,231 \\
    tokens per question (report) & 6,456 & 1,514 \\
    time per question & 23.7 s & 3.1 s \\
    \bottomrule
  \end{tabular}
\end{table}

\subsection{A degree tool}

In the large-graph pilot, node-degree questions were by far the slowest (133 seconds per answer on the MacBook,
up to 223 seconds). The time went into the first step, before any tool call: the model debated whether $G$ might be
directed and the degree therefore in-degree plus out-degree, then called \code{graph_info} to check. Its first thought
was that there is no function called degree, so \code{get_neighbors} was the way to go. Adding \code{degree}
(Table~\ref{tab:degree}) removed both the debate and the counting.

\begin{table}[ht]
  \centering
  \small
  \caption[Ablation: the degree tool]{Ablation of \code{degree}: the same 6 node-degree questions on large development
  graphs, 1 run, MacBook Air M4.}
  \label{tab:degree}
  % source: results/harness_runs/m2_pilot_large_verify, results/harness_runs/degree_tool_large_pilot, docs/learning_log.md (2026-10-07)
  \begin{tabular}{lrr}
    \toprule
    & without \code{degree} & with \code{degree} \\
    \midrule
    correct & 6/6 & 6/6 \\
    time per question & 133 s (max 223 s) & 17 s (max 23 s) \\
    output tokens (mostly thinking) & 1,988 & 264 \\
    total tokens & 3,758 & 1,582 \\
    rescued text tool calls & 4/6 & 1/6 \\
    \bottomrule
  \end{tabular}
\end{table}

\subsection{Building blocks for multi-step questions}

Two questions that no single tool answered, the farthest node from a given node and the number of nodes within two
hops, were asked on one small graph (7 nodes) and one large graph (29 nodes), first with the basic tools and then with
\code{distances_from} and \code{neighborhood} added (Table~\ref{tab:multistep}). These runs were not traced; the
numbers come from the console output recorded in the learning log.

\begin{table}[ht]
  \centering
  \small
  \caption{Multi-step probes with and without the building-block tools (1 run each, not traced).}
  \label{tab:multistep}
  % source: docs/learning_log.md (2026-10-07, multi-step questions)
  \begin{tabularx}{\textwidth}{l>{\raggedright\arraybackslash}X>{\raggedright\arraybackslash}p{4cm}}
    \toprule
    Question & Basic tools & With \code{distances_from}, \code{neighborhood} \\
    \midrule
    small, farthest & correct; 7 steps, 70 s (components, then 4 shortest paths) & correct; 2 steps, 20 s \\
    small, within 2 hops & gave up with no tool call, 87 s & correct; 2 steps, 10 s \\
    large, farthest & gave up after one shortest-path call & correct; 2 steps, 16 s \\
    large, within 2 hops & answered 26 (truth 28) after 4 of the needed neighbor calls, 151 s & correct; 2 steps, 9 s \\
    \bottomrule
  \end{tabularx}
\end{table}

With low-level tools the model would have needed about 28 calls and the merging of lists in its head, so it gave up
or did part of the work and guessed. The new tools were not written for these exact questions: they also answer ``how
many nodes at distance 3'', ``is anything unreachable'', eccentricity, and any $k$-hop question. Two side observations:
\code{cannot_answer} stopped the guessing but became an easy way out on questions the model could have answered, and
the wrong answer of 26 looked as confident as a correct one, because these question types have no verifier.

\section{M2 closing run: verifiers on and off}

\subsection{Pilot}

A pilot on the MacBook (large development graphs, all nine tasks with large questions, six questions each, one run,
verifiers on) answered all 54 questions correctly in lenient mode and 43 of 54 (79.6\%) in strict mode. In every
question the model picked the right tool on its first call. Eleven answers (20\%) were rescued text tool calls, and in
five of six cycle-check questions the verifier sent the answer back once, because the model answered ``yes'' without
the cycle. Its thinking showed why: the cycle was in the tool result, but ``the user only asked if there's a cycle, so
confirming true is sufficient''. The description of \code{submit_answer} makes the evidence sound optional.
% source: results/harness_runs/m2_pilot_large_verify, docs/learning_log.md (2026-10-07, pilot)

\subsection{Closing run}

The closing run of M2 was made on the RTX 5070 with the large development graphs: nine tasks, 15 questions each, two
runs, prompt version v2, once with verifiers (\code{m2_large_verify}) and once without (\code{m2_large_noverify}).
Without verifiers the model is still asked for the same evidence; only the checking is off. The two fixes suggested by
the pilot (saying that evidence is required, and stating in the prompt that $G$ is undirected) were not applied first.
Table~\ref{tab:m2} gives the overall numbers and Table~\ref{tab:m2-tasks} the strict accuracy per task.

\begin{table}[ht]
  \centering
  \small
  \caption{M2 closing run: large development graphs, 9 tasks $\times$ 15 questions $\times$ 2 runs, qwen3:8b, RTX 5070.}
  \label{tab:m2}
  % source: results/harness_runs/m2_large_verify, results/harness_runs/m2_large_noverify
  % (uv run python -m eval.analysis.report m2_large_verify m2_large_noverify --by-task)
  \begin{tabular}{lrr}
    \toprule
    & verifiers on & verifiers off \\
    \midrule
    answers & 270 & 270 \\
    accuracy (lenient) & 100.0\% & 100.0\% \\
    strict accuracy & 88.1\% & 90.4\% \\
    answers a verifier looked at & 44\% & 0\% \\
    answers sent back by a verifier & 18 & -- \\
    runs without an answer & 0 & 0 \\
    tokens per question & 2,237 & 2,145 \\
    time per question & 4.6 s & 4.3 s \\
    \bottomrule
  \end{tabular}
\end{table}

\begin{table}[ht]
  \centering
  \small
  \caption[M2 closing run per task]{M2 closing run per task. Lenient accuracy is 100\% for every task in both
  conditions; the table shows strict accuracy, tokens and time per question.}
  \label{tab:m2-tasks}
  % source: results/harness_runs/m2_large_verify, results/harness_runs/m2_large_noverify (report --by-task)
  \begin{tabular}{lrrrrrr}
    \toprule
    & \multicolumn{3}{c}{verifiers on} & \multicolumn{3}{c}{verifiers off} \\
    \cmidrule(lr){2-4}\cmidrule(lr){5-7}
    Task & strict & tokens & time & strict & tokens & time \\
    \midrule
    \code{connected_components_count} & 96.7\% & 2,081 & 3.4 s & 96.7\% & 2,068 & 3.2 s \\
    \code{connected_nodes} & 46.7\% & 2,368 & 5.2 s & 60.0\% & 2,135 & 4.2 s \\
    \code{cycle_check} & 93.3\% & 2,797 & 5.0 s & 96.7\% & 2,121 & 3.9 s \\
    \code{edge_count} & 100.0\% & 1,930 & 3.2 s & 100.0\% & 1,952 & 3.3 s \\
    \code{edge_existence} & 100.0\% & 1,934 & 3.5 s & 100.0\% & 1,965 & 3.1 s \\
    \code{mst} & 83.3\% & 3,128 & 11.0 s & 83.3\% & 3,163 & 11.0 s \\
    \code{node_count} & 93.3\% & 1,980 & 3.7 s & 96.7\% & 1,982 & 3.7 s \\
    \code{node_degree} & 100.0\% & 1,875 & 2.7 s & 90.0\% & 1,875 & 2.7 s \\
    \code{shortest_path} & 80.0\% & 2,042 & 3.6 s & 90.0\% & 2,046 & 3.6 s \\
    \bottomrule
  \end{tabular}
\end{table}

Four observations follow.
\begin{enumerate}
  \item The verifiers never caught a wrong answer, because there were none: the tools compute the answer and the
        model copies it correctly. On these tasks the verifiers cost a little (4\% more tokens) and bought nothing.
  \item All 18 rejections were the same: a ``yes'' to cycle check without the cycle (18 of 30 cycle-check runs), each
        fixed on the second try. This made cycle check 32\% more expensive in tokens (2,797 against 2,121) and about
        1.1 seconds slower. It is the pilot's finding again, and the open fix A (evidence required) targets it.
  \item Rescued text tool calls are the real error source: 10--12\% of answers. They concentrate in
        \code{connected_nodes}, where 16 of 30 and 12 of 30 answers were rescued, about half, and then in shortest
        path and MST: the tasks whose answers are long lists.
  \item The large dataset graphs no longer separate anything: every task is at 100\% lenient accuracy with the
        current tools. The confidence intervals in the report are degenerate (100--100\%) for the same reason.
\end{enumerate}
The differences in strict accuracy between the two conditions (for example 46.7\% against 60.0\% on
\code{connected_nodes}) come from 30 answers per cell and one pair of runs each, and should not be read as an effect
of the verifiers: the verifier only runs after an answer has been submitted, rescued or not.

\subsection{Process metrics}

Scored with the process metrics of Section~\ref{sec:process-metrics}, both M2 runs are process-perfect: tool precision,
recall and F1 of 1.00, parameter match 1.00 and exact match 100\% over all 270 answers each. On these nine tasks the
cheap model always called the one expected tool with the right arguments, and nothing else.
% source: results/harness_runs/m2_large_verify, results/harness_runs/m2_large_noverify (report --process)

\section{Large graphs: result handles}

The handle mechanism was tested on one generated graph with 10,000 nodes, 25,000 edges and 68 components
(Section~\ref{sec:handles}). After the schema fix, qwen3:8b answered the minimum-spanning-tree question in two steps,
35 seconds and 1,974 prompt plus 601 completion tokens, by submitting the handle of the tool result, and the verifier
checked all 9,932 edges. Without handles, the tool result alone would have been about 110,000 characters. This is a
single run on a single graph; size scaling is planned for M7 with the M6 generator.
% source: results/harness_runs/handles_big_graph_test, docs/learning_log.md (2026-10-07, M3 step 3)

\section{Previews of the comparison with a general-purpose harness}

Three small comparisons with Claude Code were made. None of them is the fair comparison planned for M6 (same model,
same tools, many questions): they use a different model, and the first two also different tools.

\subsection{One question, code against tools}

On one cycle-check question on a small graph, Claude Code with Sonnet 5.5, allowed only to run Python and read the
graph file, wrote one networkx script and answered correctly, using 38,660 input tokens (29,071 read from the cache and
9,585 written to it), 196 output tokens, 4.1 seconds and \$0.046 API-equivalent. Our harness with qwen3:8b answered
correctly with 1,093 input and 276 output tokens in 3.0 seconds, at zero dollar cost. The overhead, about 35 times more
input tokens, comes almost entirely from Claude Code's own system prompt and tool definitions.
% source: docs/learning_log.md (2026-10-04, code-writing agent vs tool harness)

\subsection{Large graphs, nine questions}

Table~\ref{tab:cc-large} compares Claude Code with Sonnet 5.5 (again with Python and the graph file) and our harness
with qwen3:8b on three questions each of \code{mst}, \code{connected_nodes} and \code{node_degree} on large development
graphs, one run, both graded by \code{eval/tasks.py}.

\begin{table}[ht]
  \centering
  \small
  \caption{Claude Code (Sonnet 5.5, Python and networkx) against our harness (qwen3:8b, RTX 5070): 9 questions on
  large development graphs, 1 run. Averages per question unless stated.}
  \label{tab:cc-large}
  % source: results/harness_runs/claude_code_vs_harness_large (results.jsonl, progress.log)
  \begin{tabular}{lrr}
    \toprule
    & ours (qwen3:8b, local) & Claude Code (Sonnet 5.5) \\
    \midrule
    correct & 9/9 & 9/9 \\
    input tokens per question & 1,497 & 41,925 \\
    output tokens per question & 531 & 345 \\
    time per question & 12.2 s & 6.0 s \\
    time per question, \code{node_degree} & 2.8 s & 5.3 s \\
    time per question, \code{connected_nodes} & 5.4 s & 5.4 s \\
    time per question, \code{mst} & 28.4 s & 7.2 s \\
    dollar cost, all 9 questions & \$0 & \$0.38 (API-equivalent) \\
    \bottomrule
  \end{tabular}
\end{table}

Both systems answered all nine questions. Claude Code used about 28 times more input tokens per question, mostly its
own prompt and tools (cached, but present on every call). Our harness was faster on short answers and slower on long
ones: an 8-billion-parameter model on a consumer GPU writes out a 50-edge spanning tree slowly, and the first, cold-start
question took 65.9 seconds alone (without it, our average time per question is about 6.8 seconds, as recorded in the
learning log). On these tasks networkx does the real work in both setups, so the only difference is cost.

\subsection{Claude Code with only our tools, through MCP}

In the third comparison Claude Code's built-in tools were switched off and it saw only our tools through the MCP
server, on the pilot's shortest-path question (large graph 0, from node 11 to node 21). Table~\ref{tab:mcp-demo} shows
the result. Claude Code saw exactly our twelve tools of that time and no built-in ones, and used them correctly.

\begin{table}[ht]
  \centering
  \small
  \caption{One shortest-path question: qwen3:8b in our harness (pilot, MacBook) against Sonnet in Claude Code with only
  our tools via MCP.}
  \label{tab:mcp-demo}
  % source: results/harness_runs/claude_code_mcp_demo (stream.jsonl, README.md), docs/learning_log.md (2026-10-08, M3 step 4)
  \begin{tabular}{l>{\raggedright\arraybackslash}p{4.8cm}>{\raggedright\arraybackslash}p{5.6cm}}
    \toprule
    & qwen3:8b, our harness & Sonnet, Claude Code + MCP \\
    \midrule
    answer & [11, 21], correct & [11, 21], correct \\
    calls & \code{shortest_path}, \code{submit_answer} & \code{shortest_path}, text answer \\
    time & 29.6 s (MacBook) & 3.9 s \\
    input tokens & 1,290 & about 21,500 (10,864 cache write, 10,700 cache read) \\
    output tokens & 440 (mostly thinking) & 135 \\
    dollar cost & \$0 & \$0.047 (API-equivalent) \\
    \bottomrule
  \end{tabular}
\end{table}

With identical tools, the general-purpose harness still carried about 17 times more context for the same question.
Claude Code answered in text, so grading at scale needs \code{submit_answer} served over MCP (M6). The model alias
\code{sonnet} resolved to \code{claude-sonnet-5} rather than 5.5 in this run, so the two Claude Code rows in this
chapter do not use the same model version.

\section{A multi-step task: the waypoint route}

Ten waypoint questions on small development graphs were run once on the MacBook. All ten were answered correctly, with
no verifier rejection and one rescued text call (strict accuracy 90\%), at about 2.3 graph-tool calls, 6,758 tokens
and 239 seconds per question, almost all of it thinking (one two-call question took about 10,000 thinking tokens). All
four routes that revisit a node were joined correctly, and in one run the model issued both shortest-path calls in a
single reply. Table~\ref{tab:waypoint} gives the trajectory classes. The process metrics for the same run are tool
precision 0.80, recall 0.90, F1 0.84 and parameter match 0.83, with exact match in 50\% of questions.
% source: results/harness_runs/waypoint_small_dev (report --process), docs/learning_log.md (2026-10-08, multi-step task)

\begin{table}[ht]
  \centering
  \small
  \caption{Trajectory classes of the 10 waypoint runs, before and after the description of \code{has_path} was
  corrected (same runs, re-analysed).}
  \label{tab:waypoint}
  % source: docs/learning_log.md (2026-10-08, a multi-step task and trajectory analysis)
  \begin{tabular}{lrr}
    \toprule
    Path taken & first analysis & re-analysis \\
    \midrule
    ideal & 5 & 6 \\
    extra calls & 1 & 2 \\
    alternative path & 4 & 2 \\
    wrong path / right calls but wrong answer & 0 & 0 \\
    \bottomrule
  \end{tabular}
\end{table}

In four of ten runs the model used \code{has_path} for one leg, although its description then promised only ``one
such path'', not a shortest one. It worked because the implementation returns a shortest path; the description was
corrected to say so, and the analysis now counts such a leg like a shortest-path leg. The runs themselves saw the old
description.

With \code{run_python} offered on the first three of these questions, the model never used code: it made the same two
calls and joined the paths in its head. Two of three were correct; the third ended in \code{cannot_answer}, because
the question's wording at the time (``a shortest route from A to B that passes through C'') could legitimately be read
as ``one of the shortest A-to-B routes that happens to pass through C'', of which there was none. The wording was
changed.
% source: results/harness_runs/waypoint_small_python, docs/learning_log.md (2026-10-08, M3 step 5)

\section{Code versus tools}

\subsection{Step 1: a two-call task}

The waypoint task was run on five large development questions in four setups with qwen3:8b (prompt version v3):
graph tools only; tools and \code{run_python}; code only with our tools as functions (C1); and code only with networkx
(C2). Table~\ref{tab:code-step1} shows the result.

\begin{table}[ht]
  \centering
  \small
  \caption{Code versus tools on the waypoint task: 5 large development questions, 1 run, qwen3:8b.}
  \label{tab:code-step1}
  % source: results/harness_runs/code_A_tools, code_B_tools_plus_code, code_C1_code_tools, code_C2_code_networkx (report --code)
  \begin{tabular}{lrrrrr}
    \toprule
    Setup & correct & strict & no answer & tokens / q & questions with code \\
    \midrule
    tools only & 4/5 & 20\% & 1 & 4,538 & -- \\
    tools + code & 5/5 & 20\% & 0 & 5,165 & 0\% \\
    code only, our tools (C1) & 0/5 & 0\% & 5 & 5,913 & 60\% \\
    code only, networkx (C2) & 5/5 & 20\% & 0 & 3,483 & 100\% \\
    \bottomrule
  \end{tabular}
\end{table}

A two-call task does not need code, and when code was offered next to the tools the model never used it. Forced to use
code with our tools, it failed every question: inside code the tools still returned their JSON dictionaries, so a loop
over \code{get_neighbors(4)} iterated over the dictionary's keys and error dictionaries flowed on as results. This led
to the plain-value interface described in Chapter~\ref{ch:impl}. With networkx the model succeeded with the fewest
tokens. The strict accuracy of 20\% in three setups shows that, on these runs, most answers needed a rescued text tool
call.

\subsection{Step 2: the combine tasks}

The combine tasks need many lookups and a comparison, which is where a loop in code should help. The setups of
Table~\ref{tab:setups} were run with \code{qwen3.5:9b} on 15 large development questions (five per task), one run each,
with a 16,384-token context (prompt version v6; v7 for NX, whose text is identical when no compaction happens, and none
did). An earlier round of the same experiment was invalidated by the context-window problem of
Section~\ref{sec:case-4k} and is not reported here. Tables~\ref{tab:combine} to~\ref{tab:combine-process} give the
results.

\begin{table}[ht]
  \centering
  \small
  \caption{Combine tasks, overall: qwen3.5:9b, 15 large development questions, 1 run per setup, 16k context.}
  \label{tab:combine}
  % source: results/harness_runs/combine16k_qwen35_{A,B,H,C,NX}
  % (uv run python -m eval.analysis.report combine16k_qwen35_A combine16k_qwen35_B combine16k_qwen35_H combine16k_qwen35_C combine16k_qwen35_NX --by-task --code --process)
  \begin{tabular}{lrrrrr}
    \toprule
    Setup & correct & accuracy (95\% CI) & no answer & tokens / q & time / q \\
    \midrule
    A, tools only & 2/15 & 13.3\% (0.0--33.3) & 11 & 24,978 & 95.4 s \\
    B, tools + code & 8/15 & 53.3\% (26.7--80.0) & 7 & 21,205 & 41.0 s \\
    H, tools + code + hint & 14/15 & 93.3\% (80.0--100.0) & 1 & 13,434 & 22.3 s \\
    C, code only (our tools) & 14/15 & 93.3\% (80.0--100.0) & 0 & 8,334 & 27.9 s \\
    NX, code only (networkx) & 13/15 & 86.7\% (66.7--100.0) & 2 & 11,476 & 24.8 s \\
    \bottomrule
  \end{tabular}
\end{table}

\begin{table}[ht]
  \centering
  \small
  \caption{Combine tasks, correct answers per task (out of 5).}
  \label{tab:combine-tasks}
  % source: results/harness_runs/combine16k_qwen35_{A,B,H,C,NX} (report --by-task)
  \begin{tabular}{lrrrrr}
    \toprule
    Task & A & B & H & C & NX \\
    \midrule
    \code{hop_max_degree} & 2 & 3 & 5 & 4 & 5 \\
    \code{common_neighbors_max} & 0 & 2 & 4 & 5 & 3 \\
    \code{triangle_count} & 0 & 3 & 5 & 5 & 5 \\
    \bottomrule
  \end{tabular}
\end{table}

\begin{table}[ht]
  \centering
  \small
  \caption{Combine tasks, how code was used.}
  \label{tab:combine-code}
  % source: results/harness_runs/combine16k_qwen35_{A,B,H,C,NX} (report --code)
  \begin{tabular}{lrrrrrr}
    \toprule
    Setup & used code & code + tools & code calls / q & code errors & tool calls / q & empty replies \\
    \midrule
    A & -- & -- & -- & -- & 16.0 & 2 \\
    B & 87\% & 87\% & 3.6 & 30\% & 6.5 & 1 \\
    H & 100\% & 73\% & 2.7 & 40\% & 0.9 & 0 \\
    C & 100\% & 0\% & 3.2 & 40\% & 0.0 & 1 \\
    NX & 100\% & 0\% & 4.1 & 63\% & 0.0 & 0 \\
    \bottomrule
  \end{tabular}
\end{table}

\begin{table}[ht]
  \centering
  \small
  \caption[Combine tasks, process metrics]{Combine tasks, process metrics. ``F1 right'' and ``F1 wrong'' are the mean
  F1 of runs with a correct and a wrong answer. In NX the code calls networkx, not our tool functions, so expected calls
  are rarely matched and the metrics are not meaningful for that setup.}
  \label{tab:combine-process}
  % source: results/harness_runs/combine16k_qwen35_{A,B,H,C,NX} (report --process)
  \begin{tabular}{lrrrrrrr}
    \toprule
    Setup & precision & recall & F1 & params & exact & F1 right & F1 wrong \\
    \midrule
    A & 0.11 & 1.00 & 0.19 & 1.00 & 33\% & 0.29 & 0.17 \\
    B & 0.22 & 1.00 & 0.31 & 1.00 & 20\% & 0.38 & 0.23 \\
    H & 0.53 & 1.00 & 0.64 & 1.00 & 53\% & 0.67 & 0.18 \\
    C & 0.59 & 0.97 & 0.64 & -- & 53\% & 0.64 & 0.67 \\
    NX & 0.04 & 0.23 & 0.06 & -- & 0\% & 0.05 & 0.18 \\
    \bottomrule
  \end{tabular}
\end{table}

The pattern is clear even on 15 questions. With tools only, the model made 16 graph-tool calls per question on
average and ran out of steps in 10 of 15 questions (one more ended in loop detection); it answered 2 correctly. Offered
code next to the tools, it used code in 87\% of questions and doubled its accuracy, but still ran out of steps in 7.
With the one-paragraph strategy hint it used code in every question, combined a tool call for the candidates with code
over them in 73\% of questions, and answered 14 of 15, with about half the tokens and a quarter of the time of setup A.
Code alone with our tool functions reached the same 14 of 15 at the lowest token count; code with networkx was slightly
worse (13 of 15) and had the highest share of failing code calls (63\%).

In the process metrics, recall is about 1 in every setup that uses our tools: the expected calls were always made.
Precision is what separates the setups, from 0.11 with tools only (many repeated single calls) to 0.53--0.59 with code.
Unlike the finding of the GDS Agent paper, where tool metrics barely correlate with accuracy, F1 here is higher for
right answers than for wrong ones in setups A, B and H (0.67 against 0.18 in H), but setup H has a single wrong answer,
so that contrast rests on one run.

Three caveats apply. These are 15 questions, one run each, so a single question moves accuracy by 6.7 points and the
intervals overlap for H, C and NX. Only one model was tested in valid conditions. And the answers are numbers, so none
of them could be verified: code that prints a number provides no evidence a verifier can check, which is a reason to
keep evidence-returning tools even where code is more accurate.

\section{Case study: a silent 4,096-token context}
\label{sec:case-4k}

The first round of the combine experiments (qwen3:8b and qwen3.5:9b, setups A, B, H and C) showed a puzzling failure.
qwen3.5:9b never submitted a wrong answer; its misses were runs that went silent: it found the answer (``the answer is
3'') and then sent empty replies until loop detection ended the run. A new prompt version (v6) added a dedicated nudge
for empty replies, showed the output of \code{graph_info()} in the tool description, and restored redefined tool
functions. The rerun still went silent, and the nudge did not help.

Reading the traces gave the explanation. The model's thinking stopped mid-sentence (``The output shows that''), and in
one case the model had lost the question. No call in these runs ever had a prompt longer than 4,088 tokens. Ollama
serves every model with a 4,096-token context unless asked otherwise, although qwen3.5 supports 262k tokens; a prompt
near the limit leaves no room for the reply, and nothing raises an error. In the first round, the largest prompt of
each qwen3.5:9b run was between 3,941 and 4,011 tokens; in the rerun with a 16k context, the largest prompts were
between 3,505 and 7,725 tokens, so the model did need more than 4,096 tokens.
% source: docs/learning_log.md (2026-10-08, code vs tools ... silent 4k context); max prompt tokens read from
% results/harness_runs/combine_qwen3_{A,B,H,C}, combine_qwen35_{A,B,H,C} and combine16k_qwen35_{A,B,H,C,NX} (traces.jsonl, model_call prompt_tokens)

The fix (requesting 16k tokens, logging the context window, flagging prompts above 90\% of it, and later compaction)
is described in Chapter~\ref{ch:impl}. Earlier runs were checked: they all stayed well under the limit (the M2 closing
runs, for example, never exceeded about 1.7k prompt tokens), so the M2 results and the Claude Code comparison stand.
All first-round combine runs are treated as invalid in this thesis, including those of qwen3:8b, two of whose runs
reached prompts of 4,034 and 4,077 tokens, within about 60 tokens of the limit. The qwen3:8b combine experiment has not been rerun.
\todo[inline]{Rerun the combine setups with qwen3:8b at 16k context (or with the pilot's chosen models) so that the code-versus-tools result covers more than one model.}

Why does Ollama default to so small a context? To write each new token the model attends to every earlier token, and
the keys and values of earlier tokens are kept in GPU memory, the KV cache, whose size grows linearly with the context:
$2 \times \text{layers} \times \text{KV heads} \times \text{head size} \times \text{bytes} \times \text{tokens}$. For
qwen3:8b (36 layers, 8 KV heads of size 128, 16-bit values) that is about 144~KB per token: about 0.6~GB at 4k, 2.3~GB
at 16k and 18~GB at 128k, more than the whole GPU. The cache is reserved up front, so a small default fits everywhere.
qwen3.5 uses full attention in only 8 of its 32 layers, with 4 KV heads of size 256, about 32~KB per token, so 16k
costs about 0.5~GB; the whole loaded model measured 5.8~GB.
% source: docs/learning_log.md (2026-10-08, KV cache)

The lesson is methodological: a serving setting produced behaviour that looked exactly like a model weakness. For
the thesis, the context size of every local run is now recorded in the trace and reported.

\section{Summary of findings so far}

\begin{itemize}
  \item On the existing dataset, a cheap local model with fitting tools is accurate: 100\% lenient accuracy over 270
        large-graph answers, with a process-perfect tool choice. The dataset is saturated on large graphs and cannot
        separate setups on most tasks.
  \item The remaining errors are format errors: tool calls written as text (10--12\% of answers, about half on the
        neighbor-list task) and evidence left out of a ``yes''. Reasoning errors were not observed on these tasks.
  \item A tool that fits the task cut cost by roughly an order of magnitude at equal accuracy, three times
        (\code{has_edge}, \code{degree}, \code{distances_from}/\code{neighborhood}).
  \item Verifiers caught no wrong answers on these tasks, because the tools left none; their only effect was to send
        back answers without evidence.
  \item In the previews, a general-purpose harness with a frontier model reached the same accuracy with 17--35 times
        more input tokens per question.
  \item On aggregation tasks, code over the same tools turned 2 of 15 into 14 of 15 for qwen3.5:9b, at about half the
        tokens; the strategy hint mattered.
  \item All of this is one run, few questions, the development split and mostly one model.
\end{itemize}
